\section{Introduction}
A transformer trained on many synthetic tasks and asked to continue a new sequence is often analysed against Bayes-optimal
reference predictors. Below a task-diversity threshold such models track the posterior over the finite pool of training
tasks; above it they track the Bayes predictor of the continuous task prior \citep{raventos2023,lu2025,panwar2024}. The same
memorising-versus-generalising contrast has been drawn for finite sets of discrete Markov chains \citep{gibson2026,park2025},
and a rational-strategy account describes which reference a model follows during training \citep{wurgaft2025}. In each case
the behavioural claim rests on a premise: that the candidate references predict differently on the evaluation data, by more
than the analysis can resolve.

We examine this premise for in-context forecasting of autoregressive dynamical systems, before any network is trained,
and ask of the references themselves: \emph{in which cells could a model's behaviour be classified at all?} We distinguish
three references: Bayes under the population prior from which training systems are drawn; the posterior over a finite
library of training systems (the memorising reference); and a \emph{context-only} Bayes predictor that knows the noise model
and generic structure (stability) but nothing about the population. In linear regression with a Gaussian prior, ridge
regression with the matched penalty is the continuous-prior Bayes predictor, and the usual context-only comparator is least
squares; here we make the context-only comparator itself a Bayes predictor with a deliberately uninformative prior over the
system, so that the population-specific part of the advantage can be measured.

\paragraph{Contributions.}
(i) Validated reference predictors for scalar AR(1) with known noise scale, AR(1) observed through Gaussian noise (Kalman
likelihood), AR(1) with unknown noise scale, and a bivariate VAR(1) with a stability-truncated conjugate prior, with
measured numerical error, calibration and optimality tests, and mutation tests for the last three families.
(ii) Two classical quantities, applied to the design of behavioural studies: the relative entropy between the population and
context-only priors, a budget that the summed per-forecast gaps cannot exceed (the cumulative log loss, or redundancy, of one
Bayes mixture against another), and the divergence to the nearest library system, the long-context limit of the library gap
(Section~\ref{sec:theory}). Neither bound is new; their use to predict where references can be told apart is.
(iii) Separability maps of these references at a common margin, reported cell by cell with their power limits and a
sensitivity analysis in $\delta$, which check (ii) numerically: the population window is short for a known noise scale and
narrows under observation noise; the larger gaps when the noise scale is unknown are mostly knowledge of the noise level, while
in the bivariate family their split between noise level and dynamics is order-dependent; libraries are separable from
continuous priors only when small in the
scalar families, and also when large in the bivariate family. The maps are reference computations whose specifications were
committed before their seeds were used; some were chosen after earlier results were seen, and Appendix~\ref{app:provenance}
says which. (iv) Pre-registered studies, including a null (\sSix: STOP) and a neural pilot that its own frozen rules
declared not resolvable for compute reasons (Appendices~\ref{app:history} and~\ref{sec:pilot}), together with the machinery that
kept these results honest (freeze manifests, append-only disclosure logs, an authorization gate, and a claims table that
re-derives every number in this paper from pinned commits).
